% ==========================================================================
% Capítulo: Descripción del problema
% Estado: se registra en ../STATUS.md, no aquí.
%
% Debe contener (project-context/formato-anteproyecto.md, "Descripción del
% problema"; extensión máxima 1 página):
%   - IDENTIFICACIÓN DEL PROBLEMA: síntomas y manifestaciones, con sus
%     causas y efectos jerarquizados (causas principales/secundarias,
%     efectos principales/secundarios), derivados del capítulo anterior.
%   - FORMULACIÓN DEL PROBLEMA: enunciado claro, sucinto y sin
%     ambigüedades del problema a resolver, con su pertinencia y
%     relevancia en el contexto.
% NO debe contener: la solución, los objetivos, ni tecnología concreta
% (eso va en 04-objetivos.tex y 07-metodologia.tex).
%
% El árbol de problemas correspondiente va en 99-anexos.tex.
% Hechos del proyecto: ../../project-context/documentation.md (secciones
% "Problem Description" y "Scope") y requirements.md.
% ==========================================================================
\chapter{Descripción del problema}\label{cha:descripcion-problema}

\section{Identificación del problema}

\todo{La estructura general del apartado es coherente y sigue una secuencia problema → síntomas → causas → efectos. Sin embargo, la relación causal entre estos elementos no siempre es explícita: algunas “causas secundarias” parecen más bien limitaciones o condiciones del contexto, mientras que los “efectos secundarios” provienen de problemas distintos (acceso, gobernanza y monitoreo). Conviene revisar que cada causa explique directamente un síntoma y que cada efecto se derive claramente del problema o de sus síntomas.}
\todo{"El problema central es que..." es un lenguaje muy básico. La formulación del problema resulta demasiado coloquial; convendría emplear una construcción más académica que introduzca la situación problemática.}
El problema central es que la infraestructura de GPU del IAsLab no puede aprovecharse de forma compartida, controlada y observable por la comunidad académica que la necesita. Ese problema se manifiesta en tres síntomas. Poner un modelo en operación exige un procedimiento manual, el uso de los equipos depende de credenciales que se entregan a pocas personas, y los administradores no pueden observar en tiempo real el consumo ni el estado de los nodos.

\todo{"La causa principal es que..."}
La causa principal es que la operación de la infraestructura descansa en procedimientos manuales, porque no existe una capa de \emph{software} que automatice el despliegue de modelos, reparta los recursos entre usuarios y registre su uso. A ella se suman tres causas secundarias. El proyecto de grado previo automatizó el aprovisionamiento de nodos para el entrenamiento, pero no las etapas posteriores. SAAMFI identifica a los usuarios y sus roles, sin traducir esos roles en límites de consumo. Por último, los nodos con GPU también se usan para las clases sin que exista una vista de su ocupación.

El efecto principal es la subutilización de una infraestructura en la que la universidad ya invirtió, pues la mayoría de los estudiantes, sin acceso a los equipos, recurre a servicios externos o prescinde del cómputo especializado. Los efectos secundarios recaen sobre distintos involucrados. Los administradores concentran la carga de habilitar accesos y desplegar modelos, y los estudiantes y profesores carecen de un servicio para consumir los modelos que entrenan. Si el acceso se amplía sin reglas de reparto, aumenta además el riesgo de que unos pocos usuarios o trabajos monopolicen las GPU, y de que fallas como el congelamiento de un nodo por desbordamiento de memoria pasen inadvertidas hasta afectar a otros usuarios.

\section{Formulación del problema}

\todo{La formulación es coherente con el planteamiento general, pero la pregunta incorpora directamente la solución propuesta. Revisar si se busca formular una pregunta de investigación sobre el problema o una pregunta de desarrollo tecnológico orientada a construir la plataforma. También conviene mantener consistencia entre los conceptos de roles, cupos y reglas de reparto}
La infraestructura de GPU del IAsLab no puede aprovecharse de forma compartida, controlada y observable, porque su operación depende de procedimientos manuales y ninguna herramienta automatiza el despliegue de modelos, reparte los recursos según los roles de los usuarios ni registra el estado y el consumo de los nodos. Como consecuencia, la capacidad instalada se subutiliza, la mayoría de los estudiantes queda sin acceso a ella, el ciclo de MLOps se interrumpe después del entrenamiento, y ampliar el uso sin reglas expondría las GPU a un reparto inequitativo y a fallas que nadie detecta a tiempo. De ello se deriva la pregunta del proyecto. ¿Cómo construir, sobre la infraestructura del IAsLab, una plataforma de despliegue e inferencia de modelos, reserva de cupos con reglas de reparto y observabilidad avanzada, sin comprometer la estabilidad ni la disponibilidad del \emph{hardware}? Responderla consolida en una sola plataforma las etapas hoy desarticuladas de despliegue, gobernanza y monitoreo.

% [verify: ADR-010: enlazar al árbol de problemas cuando exista en 99-anexos.tex]
